\section{Chemical Development Modelling}
\label{sec:development}

\subsection{What Development Actually Changes}

Development converts the latent image into density. Varying the development
process does not translate the characteristic curve; it \emph{reshapes} it, and
the reshaping is highly structured:

\begin{itemize}
\item \textbf{Gamma rises with development, then saturates.} More development
time develops more of the exposed crystals, steepening the curve --- but only
until essentially all exposed crystals are developed, after which further time
produces only fog. The classical gamma-versus-time family
\cite{james1977,kodakh740} is asymptotic to a
value $\gamma_\infty$ characteristic of the emulsion and developer.
\item \textbf{Fog rises without bound.} Unexposed crystals develop
spontaneously at a slow rate, raising $\Dmin$ indefinitely. This is why heavy
pushing produces flat, milky shadows: the shadow density is fog, and fog
carries no image.
\item \textbf{Speed gain is much smaller than the exposure deficit.} A
two-stop push does not recover two stops. It recovers, typically, one third to
one half of a stop of true speed, because the toe crystals that received too
few photons never formed a developable latent image and no amount of
development creates one. The rest of the apparent gain is contrast increase.
\end{itemize}

The last point is the one that consumer applications universally miss, and it
is the reason \EM{}'s push processing looks correct: pushed images have
\emph{empty, contrasty shadows}, not brighter shadows.

\subsection{Development Activity}

All process variables --- time, temperature, agitation, developer strength ---
are collapsed into a single scalar \emph{development activity} $A$, normalized
so that $A = 1$ is the manufacturer's normal process. Time and temperature
combine through an Arrhenius relation on the reaction rate:

\begin{equation}
A \;=\; \frac{t}{t_0}\,\exp\!\left[\frac{E_a}{R}\left(\frac{1}{T_0} - \frac{1}{T}\right)\right]\, \eta(a_g)\, \frac{[\mathrm{dev}]}{[\mathrm{dev}]_0},
\label{eq:activity}
\end{equation}

where $t$ is development time, $T$ is absolute temperature in kelvin, $E_a$ is
the apparent activation energy of the development reaction, $R$ is the gas
constant, $\eta(a_g)$ is an agitation efficiency factor, and the final factor
is relative developer concentration. For standard color negative chemistry
$E_a/R \approx 8.4 \times 10^3$\,K, which yields the familiar rule of thumb that
a $\pm 0.5$\textdegree C error at 38\textdegree C is roughly a $\pm 4.5\%$ change
in effective time. Reproducing that sensitivity is what makes the temperature
control feel like a real process control rather than a slider.

Agitation efficiency uses a saturating form,

\begin{equation}
\eta(a_g) = \eta_\infty - (\eta_\infty - \eta_0)\,e^{-a_g/a_c},
\end{equation}

normalized so that $\eta(1) = 1$ at the recommended agitation scheme. Reduced
agitation ($a_g < 1$) both lowers activity and, importantly, \emph{increases}
the interlayer adjacency effect of Section~\ref{sec:interlayer}, because
exhausted developer and released inhibitor are not swept away from the emulsion
surface. This coupling is modelled explicitly and is the mechanism behind
stand-development's characteristic edge effects.

\subsection{Parameter Modulation}

Given $A$ and a push/pull setting of $n$ stops (positive is push), the stock's
nominal parameters $\theta^{(0)}$ are modulated as follows. Push and pull are
expressed through $A$ --- a push of $n$ stops corresponds to
$A = A_0 \cdot \rho^{\,n}$ with $\rho \approx 1.35$ per stop for color negative
chemistry --- so there is a single mechanism, not two.

\paragraph{Gamma} follows the saturating development curve:

\begin{equation}
\gamma(A) \;=\; \gamma_\infty\left(1 - e^{-A/A_\gamma}\right),
\label{eq:gammatime}
\end{equation}

with $A_\gamma$ chosen per stock so that $\gamma(1) = \gamma^{(0)}$, i.e.

\begin{equation}
A_\gamma = \left[\ln\!\left(\frac{\gamma_\infty}{\gamma_\infty - \gamma^{(0)}}\right)\right]^{-1}.
\end{equation}

Typical values are $\gamma_\infty \approx 1.6\,\gamma^{(0)}$, so that gamma can
rise substantially but not without limit. This single equation is why extreme
pushes stop adding contrast --- the correct behavior, obtained for free.

\paragraph{Fog} rises monotonically and without saturation:

\begin{equation}
\Dmin(A) \;=\; \Dmin^{(0)} + \phi_0\left(A^{\,\beta_f} - 1\right),
\label{eq:fog}
\end{equation}

with $\beta_f \approx 2.1$ and $\phi_0 \approx 0.045$ for color negative
stocks. The per-channel $\phi_0$ is not equal --- the blue-sensitive layer fogs
fastest --- which produces the characteristic cool-shadow crossover of
heavily pushed color film.

\paragraph{Speed} shifts by a fraction of the nominal push:

\begin{equation}
x_0(A) \;=\; x_0^{(0)} \;-\; \alpha_s \log_{10} A,
\label{eq:speedshift}
\end{equation}

with the \emph{speed recovery coefficient} $\alpha_s \in [0.25, 0.45]$,
substantially less than the $1.0$ that would represent full speed recovery.
Table~\ref{tab:pushresp} gives the resulting behavior, which matches published
push-processing data closely.

\paragraph{Toe softness} increases with development, because heavy development
extends the toe region as marginally-exposed crystals begin to develop:

\begin{equation}
\kappa_t(A) = \kappa_t^{(0)}\left(1 + \tau_t \ln A\right), \quad \tau_t \approx 0.30.
\end{equation}

\paragraph{Density range} rises modestly, since the shoulder is set primarily
by silver and coupler availability rather than by development:

\begin{equation}
\Delta D(A) = \Delta D^{(0)}\left(1 + \tau_D \ln A\right), \quad \tau_D \approx 0.12.
\end{equation}

\begin{table}[htbp]
\caption{Modelled Push/Pull Response, Typical Color Negative}
\label{tab:pushresp}
\begin{center}
\renewcommand{\arraystretch}{1.15}
\begin{tabular}{@{}lccccc@{}}
\toprule
\textbf{Setting} & $A$ & $\gamma$ & $\Delta\Dmin$ & \textbf{True} & \textbf{CI} \\
 & & & & \textbf{speed} & \\
\midrule
Pull 2  & $0.55$ & $0.44$ & $-0.02$ & $-0.24$\,EV & $0.41$ \\
Pull 1  & $0.74$ & $0.53$ & $-0.01$ & $-0.13$\,EV & $0.49$ \\
Normal  & $1.00$ & $0.61$ & $\phantom{-}0.00$ & $\phantom{-}0.00$\,EV & $0.58$ \\
Push 1  & $1.35$ & $0.69$ & $+0.05$ & $+0.15$\,EV & $0.66$ \\
Push 2  & $1.82$ & $0.77$ & $+0.13$ & $+0.29$\,EV & $0.73$ \\
Push 3  & $2.46$ & $0.83$ & $+0.28$ & $+0.42$\,EV & $0.79$ \\
\bottomrule
\multicolumn{6}{@{}p{0.95\columnwidth}@{}}{\footnotesize A three-stop push yields $0.42$\,EV of genuine speed. The remaining apparent gain is contrast, and the cost is $+0.28$ of fog in the shadows.}
\end{tabular}
\end{center}
\end{table}

\subsection{Cross-Processing}

Cross-processing --- developing color negative film in reversal chemistry
(C-41 in E-6) or the reverse --- is modelled as a substitution of the entire
modulation parameter set rather than as an extra effect. The mechanism is a
\emph{chemistry profile} that supplies $(\gamma_\infty, \phi_0, \beta_f,
\alpha_s, \tau_t, \tau_D)$ and a per-channel activity scaling
$\mathbf{A}_{\mathrm{chem}}$. Because the layers of a stock designed for one
chemistry respond with wildly different efficiency to another, the per-channel
activity scaling is far from unity --- for E-6 chemistry applied to a C-41
stock, approximately $(1.8, 1.1, 0.6)$ --- which is precisely what produces
cross-processing's extreme contrast and violent color crossover. Again, the
effect emerges from the model rather than being authored.

\begin{table}[htbp]
\caption{Chemistry Profiles}
\label{tab:chemistry}
\begin{center}
\renewcommand{\arraystretch}{1.15}
\begin{tabular}{@{}lccccc@{}}
\toprule
\textbf{Process} & $\gamma_\infty/\gamma^{(0)}$ & $\phi_0$ & $\beta_f$ & $\alpha_s$ & $\rho$ \\
\midrule
C-41 (color neg.)   & $1.60$ & $0.045$ & $2.1$ & $0.35$ & $1.35$ \\
E-6 (reversal)      & $1.35$ & $0.020$ & $1.7$ & $0.28$ & $1.30$ \\
ECN-2 (motion pic.) & $1.55$ & $0.030$ & $2.0$ & $0.33$ & $1.32$ \\
B\&W, general       & $2.10$ & $0.035$ & $2.4$ & $0.42$ & $1.45$ \\
B\&W, compensating  & $1.45$ & $0.025$ & $1.9$ & $0.38$ & $1.25$ \\
\bottomrule
\end{tabular}
\end{center}
\end{table}

\subsection{Where This Executes}

All of Section~\ref{sec:development} is host-side arithmetic. The development
model consumes the stock profile and the user's process parameters and produces
a modified $\theta$, which is uploaded as a small constant buffer. Nothing in
this section costs a GPU instruction, which is worth stating explicitly:
substantial perceived sophistication in the product is free at render time
because it modulates parameters rather than pixels.

This has an important corollary for the UI. Development controls can update at
unlimited rate without any incremental render cost beyond one re-render, and
they invalidate the baked pointwise LUT (Section~\ref{sec:lutbake}) but nothing
spatial. The dependency tracking in Section~\ref{sec:rendergraph} encodes this.
